El prototipado digital y la simulación computacional se han consolidado como herramientas fundamentales en la industria automotriz y el automovilismo deportivo, permitiendo evaluar diseños en etapas tempranas y entrenar pilotos en entornos interactivos seguros (Driver-in-the-Loop). Sin embargo, el sector se encuentra dominado por plataformas comerciales cerradas, lo que evidencia una carencia de alternativas de código abierto basadas en técnicas de modelado modernas. Para abordar esta problemática, el presente trabajo propone el desarrollo de un simulador de dinámica vehicular interactivo, estrictamente modular y capaz de ejecutarse en tiempo real. El sistema propuesto se articula en tres componentes interoperables: un núcleo físico-matemático desarrollado en Modelica, que aprovecha su paradigma de lenguaje orientado a objetos para modelar la cinemática del chasis y la suspensión, una biblioteca de comunicación de baja latencia en C++ para la integración con hardware de simulación mediante las API SDL y DirectInput, y un entorno de visualización 3D sincrónico implementado en Unity. El objetivo es lograr una arquitectura de software desacoplada que facilite el intercambio de subsistemas mecánicos, validando la estabilidad numérica y el rendimiento del lazo cerrado conductor-simulador. De este modo, se busca sentar las bases y documentar una herramienta extensible que permita el acceso a la simulación vehicular de alta fidelidad para la investigación y el desarrollo independiente.